\section{Next Token Prediction 的缺陷：压缩不等于计算精度}

前面的问题发生在视觉，本章转向语言模型，因为张祥雨认为语言模型里也出现了类似的瓶颈。训练大模型时，他观察到一个怪现象：模型越大，通用对话、情商、知识量更强；但数学等局部推理能力可能先上升、后平缓，再继续扩大反而下降。一个线索是，大模型更倾向于“跳步”，不像小模型那样老老实实算。

他的解释指向 next token prediction。这个目标本质上是在建模数据分布，也可理解为压缩：模型越能预测下一个 token，就越能压缩语料。但压缩人类语料的分布，不等于在数学、逻辑、几何和物理任务上保持计算精度。人类语料中有大量跳步、直觉和省略，大模型压缩得越好，可能越学到这些捷径。

\[
L_{\mathrm{NTP}}=-\sum_{t}\log p_{\theta}(x_t\mid x_{<t})
\]

其中，$x_t$ 是第 $t$ 个 token；$x_{<t}$ 是此前上下文；$p_{\theta}$ 是模型给出的下一个 token 概率；$L_{\mathrm{NTP}}$ 越低，说明模型越能拟合训练分布。但数学题要求的是步骤正确和结果正确，而不是输出分布看起来像互联网解答。

\begin{figure}[H]
\centering
\includegraphics[width=0.96\textwidth]{next-token-compression-defect.png}
\caption{Next Token Prediction 的缺陷：压缩分布不等于计算精度，越大会越倾向跳步。自制概念图，依据 00:41:11--00:50:48 对谈内容整理。}
\end{figure}

\begin{knowledgebox}{读图：shortcut 是分布学习的副作用}
图中从 Data Distribution 到 Compression，再到 Shortcut、Math 和 RL。关键是 shortcut：如果训练分布里有大量省略步骤，模型会学到“像人类答案”的压缩，而不是“每一步都计算正确”的程序。
\end{knowledgebox}

\subsection{特征坍缩：高概率特征与精确约束的冲突}

本节解释访谈里的“特征坍缩”。生成模型会偏向保留数据分布里最容易压缩、最常见、最稳定的特征；但复杂推理关心的往往是少数关键步骤。数学里一步错全错，图像里一个手指、透视或物理接触错了，整体就失真。分布相似不能替代约束满足。

\begin{figure}[H]
\centering
\includegraphics[width=0.94\textwidth]{feature-collapse.png}
\caption{特征坍缩现象：生成模型会保留高概率特征，却丢失计算和约束精度。自制概念图，依据 00:45:42--00:50:48 对谈内容整理。}
\end{figure}

\begin{knowledgebox}{读图：像训练语料，不等于算对}
左侧“分布接近”解释为什么模型语言流畅、知识丰富；右侧“精度不足”解释为什么数学、几何和物理仍会错。多模态推理需要的是精确约束，而不是只让输出更像训练样本。
\end{knowledgebox}

\begin{warningbox}{不要把压缩率当作智能上限}
更好的 next token loss 往往提升通用能力，但它不是所有智能任务的充分目标。需要精确搜索、长期规划、外部验证和动作反馈的任务，必须引入额外目标和过程。
\end{warningbox}

\subsection{本章小结}

next token prediction 是大模型起飞的基础，但它也解释了为什么更大模型可能更会跳步。数学、视觉可控生成和多模态推理，都需要超出分布压缩的目标函数。

